\section{A square gap that removes the radical weight}
\label{sec:gap}

The parent uses $9B_\tau$ in place of $B_\tau$, so its signed radical
relation contains $3\tau$. Deleting that weight requires more than checking
the intended integral witnesses: a rational signed sum of irrational
roots could otherwise introduce extra solutions. The following argument
excludes every such cancellation before any source semantics is used.

\begin{lemma}[A square gap]
\label{lem:gap}
Let $x\ge2$ and $t\ge1$ be integers. Put
\[
\Delta=t(x+1)(t(x+1)+2),\qquad v=t^2x,
\quad D=xt(x+1)(t(x+1)+2)(t^2x+1).
\]
Then $\Delta v(v+1)=t^2D$ is not a rational square.
In particular $\Delta I^2=v(v+1)k^2$ has no solution with nonzero
integers $I,k$.
\end{lemma}

\begin{proof}
Set $m=\lfloor x/2\rfloor$ and $h=t^2x(x+1)+tx+m$.
For $x=2m$, direct expansion gives
\[
D-h^2=m[2t^2+4(m+1)t-m]>0,
\tag{23}
\]
\[
(h+1)^2-D=2m(4m+1)t^2-4m^2t+(m+1)^2>0.
\tag{24}
\]
Here $m\ge1$. The first is positive already after replacing $t$ by 1.
In the second, $t^2\ge t$ gives the positive lower bound
$(4m^2+2m)t+(m+1)^2$.
For $x=2m+1$, where $m\ge1$, the exact identities are
\[
D-h^2=t^2x(x+2)+2tx(m+2)-m^2>0,
\tag{25}
\]
\[
(h+1)^2-D=[tx-(m+1)]^2>0.
\tag{26}
\]
The first is positive at $t=1$ and increases thereafter. In the second,
$tx\ge2m+1>m+1$. Thus $h^2<D<(h+1)^2$ in both parity classes.
An integer that is a rational square is an integer square, by reduction
of its rational square root to lowest terms.
Finally the prohibited equation would imply $D=(\Delta I/(tk))^2$.
The excluded point $x=t=1$ gives $D=16$; the domain is essential.
\end{proof}

\begin{lemma}[Rationality of three separated radicals]
\label{lem:rational}
Let $a,b,c$ be positive, pairwise distinct rational numbers with
\[
(a+b-c)^2\ne4ab.
\tag{27}
\]
If $u^2=a$, $v^2=b$, $w^2=c$ and $u+v+w$ is rational, then
$u,v,w$ are rational.
\end{lemma}

\begin{proof}
Write $s=u+v+w$. If $s=0$, squaring $w=-u-v$ twice contradicts
(27), so $s\ne0$. For $K_a=s^2+a-b-c$, squaring
$2vw=K_a-2su$ gives
\[
4sK_a u=K_a^2+4s^2a-4bc.
\tag{28}
\]
At least one of $K_a,K_b,K_c$ is nonzero, since otherwise subtracting
them would give $a=b=c$. The corresponding identity makes one root
rational. The sum of the other two is rational and nonzero, since their
squares differ. If that sum is $r=v+w$, then
$v=(r^2+b-c)/(2r)$ and $w=r-v$.
\end{proof}

\begin{theorem}[Unweighted separation]
\label{thm:separation}
Use $\Delta,v$ from Lemma~\ref{lem:gap}. Suppose an integer $\mu>3$
divides $t$, integers $I,k$ satisfy $I\ge k\ge1$, and $V$ is a
nonzero integer. Put
\[
d_\tau=1+v(v+1)k^2,\quad d_f=1+\Delta I^2,\quad
Q=\Delta^2I^2,\quad d_y=\frac{QV^2-1}{Q-1}.
\tag{29}
\]
Every rational signed sum of $\sqrt{d_\tau},\sqrt{d_f},\sqrt{d_y}$
has all three radicals rational.
\end{theorem}

\begin{proof}
We have $\Delta>v+1>1$, so $Q>1$ and all three radicands are positive.
Lemma~\ref{lem:gap} excludes $d_\tau=d_f$.
If $d_\tau=d_y$, then $Q(d_\tau-V^2)=d_\tau-1$,
but its right side is strictly between 0 and $Q$, since
$v(v+1)k^2<\Delta^2I^2$.
If $d_f=d_y$, rearrangement gives $d_f-V^2=1/\Delta$,
contradicting integrality. Thus they are pairwise distinct.

For the remaining triangle obstruction, put $\delta=Q-1$ and
$D_{\rm n}=QV^2-1$. Equality in (27) would imply
\[
[\delta(d_\tau+d_f)-D_{\rm n}]^2-4\delta^2d_\tau d_f=0.
\tag{30}
\]
Modulo $\mu$ we have $\Delta=v=Q=0$,
$\delta=D_{\rm n}=-1$, and $d_\tau=d_f=1$.
The left side is $-3$ modulo $\mu$, impossible for $\mu>3$.
Lemma~\ref{lem:rational} applies, with arbitrary signs on the roots.
\end{proof}

\subsection{Application on the full positive input domain}

In (16), $X\ge32$, $t=\mu Y$, $\mu>3$, $c=kY+\eta>k$ and
$I=ic^2\ge c^2>k$. Also $D_f>2c$ and $b$ is integral.
If $b\le0$, then $V=D_fb-c\le-c$; if $b\ge1$, then
$V\ge D_f-c>c$. Thus $V\ne0$ before the polynomial equation is used.
Every hypothesis of Theorem~\ref{thm:separation} is unconditional on the
supplied positive domain.

The octic identity
\[
F_8(s;a,b,c)=[(s+c-a-b)^2-4(sc+ab)]^2-64scab
\tag{31}
\]
factors, with $s=x^2$, into the eight signed linear radical factors.
Since $M,\delta>0$, equation (21) implies
\[
\frac ZM=\epsilon_1\tau+\epsilon_2 f+\epsilon_3 y,\qquad
\tau^2=D_\tau,\quad f^2=D_f,\quad y^2=D_{\rm n}/\delta,
\tag{32}
\]
with positive roots.
Theorem~\ref{thm:separation} makes all three roots rational.
The first two are integers because their squares are integers; the
retained adjacent-Pell descent lemma makes the third integral under the
unchanged $Q=\Delta^2I^2$ and integer-$V$ hypotheses \cite{zero}.

\section{The auxiliary Pell index supplies its own transport bound}
\label{sec:transport}

Removing the weight alone gives 98 operations at degree 368.
We now show that the already encoded auxiliary Pell solution grows enough
to remove the separate $HI$ term in the transport majorant.
This is the additional reduction to (21).

For $A\ge2$, write
\[
\chi_A(m)+\psi_A(m)\sqrt{A^2-1}
=(A+\sqrt{A^2-1})^m.
\]
The unchanged strong/auxiliary folding argument, after the integral roots
are recovered, gives $c=\psi_A(p)$ for an odd positive $p$.
The parity does not require the later equality $p=R$:
the auxiliary index is odd, and, with $m$ the strong Pell index, both
reduction modulo $2m$ and reflection
from an index $u$ to $2m-u$ preserve oddness.
Since $c\ge3$, $p=1$ is excluded, hence
\[
c\ge\psi_A(3)=4A^2-1=4\Delta+3.
\tag{33}
\]

Before transport, the retained first-index growth and mask estimate give
\[
c>8X\mu^2Y^2,\qquad R<c\alpha/4.
\tag{34}
\]
These bounds do not use $C>0$, the main index equality, or the sign of
$B_t$. In particular
\[
b=c(T+\alpha)-R>c(T+3\alpha/4)>c.
\tag{35}
\]

\begin{lemma}[Auxiliary index at least five]
\label{lem:auxfive}
Under the recovered kernel equations and (33)--(35),
\[
b>8Q>2HI.
\tag{36}
\]
\end{lemma}

\begin{proof}
Set $S_0=\Delta I$, so $Q=S_0^2$.
The auxiliary equation is
$S_0^2V^2-(S_0^2-1)y^2=1$.
Here $V>0$. Pell classification gives
$S_0V=\chi_{S_0}(\nu)$ at a positive index $\nu$.
Even indices have $\chi_{S_0}(\nu)\equiv\pm1\pmod{S_0}$,
whereas $S_0$ divides $S_0V$, so $\nu$ is odd.

By (33)--(35),
\[
V=D_fb-c>\Delta I^2c>4\Delta^2I^2=4Q.
\]
The possibilities $\nu=1$ and $\nu=3$ would give respectively
$V=1$ and $V=4Q-3$, and are impossible.
Therefore $\nu\ge5$, and monotonicity gives
\[
V\ge16Q^2-20Q+5>8Q^2.
\tag{37}
\]
The last inequality holds for $Q\ge3$, far below the actual lower bound.
Because $D_f=\Delta I^2+1<Q$,
$b=(V+c)/D_f>8Q$.
Finally $\Delta-H=(A-2)^2>0$ gives $8Q>2HI$.
\end{proof}

\begin{lemma}[Quarter bounds for the two error terms]
\label{lem:quarter}
With $u=y-V$, the recovered positive roots satisfy
\[
0<u<b/4,\qquad e/\mu^2<b/4.
\tag{38}
\]
\end{lemma}

\begin{proof}
The auxiliary equation gives
\[
u=\frac{V^2-1}{(Q-1)(y+V)}
<\frac{V}{2(Q-1)}<\frac{D_fb}{2(Q-1)}.
\]
Here $y>V>1$, and
$Q-1-2D_f=\Delta(\Delta-2)I^2-3>0$.
Thus $u<b/4$.
Also $b-e=c(T+\alpha)-k+1>0$, so $e<b$.
Since $b>0$ and $\mu^2>4$, $e/\mu^2<b/4$.
This upper bound also holds when $e$ is negative; no premature
divisibility or positivity of $e$ is assumed.
\end{proof}

\subsection{Soundness of the shorter carrier}

The retained bound $3\tau+f<HI$ implies $\tau+f<HI<b/2$.
The normalized center from (18) is
\[
\frac ZM=n+V+\frac{K_{\rm pos}}H+\frac Y{q^3}
-\frac{WB_t}{q-1}-\frac e{\mu^2}.
\tag{39}
\]
Suppose $B_t\le-1$.
With the new $W=1+(q-1)b$, $WB_t/(q-1)<-b$.
The parent growth estimate already gives $K_{\rm pos}>0$.
Using (32), (36) and (38),
\[
\begin{split}
n
&=\epsilon_1\tau+\epsilon_2f+(\epsilon_3y-V)
-K_{\rm pos}/H-Y/q^3+WB_t/(q-1)+e/\mu^2\\
&<b/2+b/4+b/4-b=0,
\end{split}
\]
contrary to $n>0$. Thus $B_t\ge0$ and hence $C>0$.

Everything after this transport step is unchanged.
In the original order the proof recovers $3q<R<q^4-2$, $p=R$, parity,
and then $\mu^2\mid e$. The first index is $(R+1)/2$.
The ratio estimate and whole-product projection force $X=2^R$;
no noninvertible factor is cancelled modulo $H$.
The remaining transport and half-binomial conditions reconstruct the
same unscaled native source and a finite syntactic refutation.
The new $W$ retains residue $1$ modulo $q-1$.

\subsection{Completeness}

Use the parent's fresh canonical zero-offset witnesses from a finite
refutation. Retain every coordinate except $n$. At these witnesses
\[
s_0=Y/q^3\ge1,\quad t_0=B_t/(q-1)\ge0,\quad
\rho=(k-R-1)/\mu^2\ge1,\quad
u=y-V>0,\quad\kappa=K_{\rm pos}/H>0
\]
are integers. With the new $W$, choose
\[
n=Wt_0+\rho+\tau+f+u-s_0-\kappa.
\tag{40}
\]
The parent gives $\tau>k>Y\ge s_0$ and $f>\kappa$.
Also $W>0$ because $b>c>0$, so $n>0$.
Equation (39) becomes $Z/M=\tau+f+y$, and a norm factor vanishes.
This is canonical reconstruction, not an asserted positive pointwise
inverse from every parent tuple.

\begin{corollary}
At the coherent coefficients of Appendix~\ref{app:constants},
$P_{96}$ has a positive integer zero exactly when the pinned ZFC calculus
has a finite refutation. Therefore ordinary mathematics formalizable in
ZFC gives (1) with $D=\Pos$ and $n=10$.
\end{corollary}

\section{Exact degrees and the direct-base companion}
\label{sec:degree}

The degree must be calculated for the entire norm. Following only the
centered difference can miss a different leading term. Write
\[
s=\delta Z^2,\quad a=\delta M^2D_\tau,\quad
b_*=\delta M^2D_f,\quad c_*=M^2D_{\rm n},\quad d=s-c_*.
\]
Then
\[
E=d^2-2(s+c_*)(a+b_*)+(a-b_*)^2.
\tag{41}
\]
For the new scaled carrier the relevant degrees are
\[
\deg a=74,\quad\deg b_*=76,\quad
\deg s=\deg c_*=104,\quad\deg d=86.
\]
The term $d^2$ has degree 172, but $-2(s+c_*)b_*$ has degree 180
and dominates it. Since $s$ and $c_*$ have the same leading component,
\[
E_0=-4(c_*)_0(b_*)_0,\qquad
(P_{96})_0=16(c_*)_0^2(b_*)_0^2.
\tag{42}
\]
The subtracted norm term has degree 358, below 360.
Thus the exact degree is \textbf{360}.
Its full leading component is
\[
2^{20}q_0^{184}w^{36}i^{20}Y^{76}
(\eta+\zeta)^{40}(q_0-F)^4,\qquad q_0=\beta J.
\tag{43}
\]
For example,
\[
[J^{184}w^{36}i^{20}Y^{76}\eta^{40}F^4]P_{96}
=2^{20}\beta^{184}>0.
\tag{44}
\]

\subsection{The direct-base equation of degree 332}

Replace $w$ by a positive input $x$ and set $X=q+x$.
Keep the same zero-offset main parameter and radicands. Use
\[
B=(qK_0+X)C-qF,\quad W=1+q(q-1)b,\quad
P_0=q^4(q-1),\quad M=P_0\mu^2H,
\]
\[
\begin{split}
Z={}&P_0\mu^2[H(n+V)+K_{\rm pos}]\\
&+H[\mu^2(q(q-1)Y-q^3WB-(q-1)X)-P_0e].
\end{split}
\tag{45}
\]
Evaluate the same unweighted norm (19)--(21).
The square gap and auxiliary index argument use $X\ge32$, not $q\mid X$,
so apply before the direct carrier recovers that divisibility.
The retained inequality $3\tau+f+X/q^4<HI<b/2$ supplies the corresponding
transport estimate.

For canonical completeness the direct parent has the exact map
\[
x=X-q,\qquad
n_{\rm direct}=n_{\rm scaled}+(q-1)^2bt_0+X/q^4>0.
\tag{46}
\]
Its numerator and denominator scale by $q$ and the norm output by $q^8$.
The late-divisibility proof and original source reconstruction remain valid
\cite{direct}.

The direct evaluator has 102 operations.
Its norm degrees are $(68,70,96,96,80)$, so the degree is 332.
The leading component is
\[
2^{20}q_0^{156}X_0^{36}i^{20}Y^{76}
(\eta+\zeta)^{40}(q_0-F)^4,\qquad X_0=q_0+x.
\tag{47}
\]
The coefficient of $J^{156}x^{36}i^{20}Y^{76}\eta^{40}F^4$ is
$2^{20}\beta^{156}>0$.
This companion trades six more operations for 28 fewer degrees.

\subsection{Retained intermediate equations}

The packet retains the weight-only reductions, before shortening the
transport carrier: 98 operations at degree 368 and 104 at degree 340.
They isolate the new radical-separation theorem.
The final 96/360 and 102/332 endpoints improve both degree and operations.

For the weight-only step the exact identity
\[
P_9-P_1
=16a(5a-s-b_*-c_*)(E_9+E_1)-512sc_*ab_*
\tag{48}
\]
has nonzero differences of exact degrees 362 and 334 respectively,
where $E_\lambda=d^2-2(s+c_*)(\lambda a+b_*)+(\lambda a-b_*)^2$.
Even this first step changes the polynomial, unlike the integer norm
rewrites.

